\documentclass[11pt,a4paper]{article}
% Common preamble for the hint sheets (\input from each file)
\usepackage[utf8]{inputenc}
\usepackage[T1]{fontenc}
\usepackage{lmodern}
\usepackage{amsmath,amssymb,amsthm,mathtools}
\usepackage[a4paper,margin=2.3cm]{geometry}
\usepackage{enumitem}
\usepackage{xcolor}
\usepackage[hidelinks]{hyperref}
\renewcommand{\leq}{\leqslant}
\renewcommand{\le}{\leqslant}
\renewcommand{\geq}{\geqslant}
\renewcommand{\ge}{\geqslant}
\setlength{\parindent}{0pt}
\setlength{\parskip}{0.35em}
% \hintstitle{Sheet title}{Class work / Homework}
\newcommand{\hintstitle}[2]{\begin{center}{\LARGE\bfseries Hints: #1\par}\vspace{0.5em}{\large #2\par}\end{center}
\noindent\textit{The hints to each problem are cumulative. Read one, try again, and open the next only if you are still stuck.}\par\medskip\hrule\medskip}
\newcommand{\problem}[1]{\par\bigskip{\large\bfseries Problem #1}\par\nopagebreak\smallskip}
\newcommand{\subpart}[1]{\par\smallskip\textbf{(#1)}\par\nopagebreak}
\newcommand{\hint}[1]{\par\smallskip\noindent\textcolor{black!70}{\textbf{Hint #1.}}\enspace}
\newcommand{\note}[1]{\par\smallskip\noindent\textit{Note.} #1\par}
\providecommand{\lt}{<}
\providecommand{\gt}{>}
\providecommand{\tightlist}{\setlength{\itemsep}{0pt}\setlength{\parskip}{0pt}}
\emergencystretch=3em

\begin{document}
\hintstitle{Linear algebra -- 1. Determinants and invertibility}{Class work}

\problem{1}
\hint{1}
Regard \(\det A\) as a polynomial \(F\) in the \(n^2\) entries \(a_{ij}\). By the chain rule, \(\frac{d}{dt}\det A(t)=\sum_{i,j}\frac{\partial F}{\partial a_{ij}}\,a_{ij}'(t)\), so it suffices to compute each partial derivative.
\hint{2}
Expand along row \(i\): \(\det A=\sum_k a_{ik}(-1)^{i+k}M_{ik}\), and the cofactors \((-1)^{i+k}M_{ik}\) do not involve any entry of row \(i\). Hence \(\partial F/\partial a_{ij}=(-1)^{i+j}M_{ij}=(\operatorname{adj}A)_{ji}\), and \(\sum_{i,j}(\operatorname{adj}A)_{ji}a'_{ij}=\operatorname{tr}(\operatorname{adj}A\cdot A')\).
\hint{3}
For \(\det e^{A}\): put \(f(t)=\det e^{tA}\). Jacobi gives \(f'(t)=\operatorname{tr}(\operatorname{adj}(e^{tA})e^{tA}A)=f(t)\operatorname{tr}A\) (since \(e^{tA}\) is invertible, \(\operatorname{adj}(e^{tA})=f(t)e^{-tA}\)), and \(f(0)=1\). For \(\operatorname{adj}(\operatorname{adj}A)\): first treat invertible \(A\), where \(B=\operatorname{adj}A=(\det A)A^{-1}\) and \(\det B=(\det A)^{n-1}\); plug into \(\operatorname{adj}B\cdot B=(\det B)E\). Then extend to all \(A\) because both sides are polynomials in the entries and invertible matrices are dense (or: polynomial identity holding on a Zariski-dense set).

\problem{2}
\hint{1}
The matrix is ``staircase''-shaped: column \(j\) equals column \(j-1\) except in rows \(i\ge j\). Subtract column \(j-1\) from column \(j\), for \(j=n,n-1,\dots,2\).
\hint{2}
After these operations the matrix is lower triangular with diagonal entries \(1\) and \(\frac1k-\frac1{k-1}=-\frac{1}{k(k-1)}\) for \(k=2,\dots,n\). Multiply them: the answer is \(\det A=\dfrac{(-1)^{n-1}}{(n-1)!\,n!}\).

\problem{3}
\hint{1}
The matrix is \(D+uu^T\) with \(D=\operatorname{diag}(2,3,\dots,n)\) and \(u=(1,\dots,1)^T\). Use the rank-one update formula (Theorem 2).
\hint{2}
\(\det D=n!\) and \(u^TD^{-1}u=\tfrac12+\dots+\tfrac1n\), so \(D_n=n!\,\bigl(1+\tfrac12+\dots+\tfrac1n\bigr)\). So \(D_n/n!=H_n\), the harmonic number, which is unbounded: the answer is no.

\problem{4}
\hint{1}
\(\cos((j+k)\theta)=\cos j\theta\cos k\theta-\sin j\theta\sin k\theta\), so \(A\) has rank at most \(2\): \(A=UV^T\) with \(U,V\) of size \(n\times 2\). Use \(\det(E+UV^T)=\det(E_2+V^TU)\).
\hint{2}
Take \(U\) with columns \((\cos k\theta)_k,(\sin k\theta)_k\) and \(V\) with columns \((\cos k\theta)_k,(-\sin k\theta)_k\). Then \(V^TU\) is the \(2\times2\) matrix of sums \(\sum\cos^2k\theta\), \(\sum\cos k\theta\sin k\theta\), \(\sum\sin^2 k\theta\) (with signs).
\hint{3}
With \(\theta=2\pi/n\) and \(n\ge3\) one has \(\sum_{k=1}^n\cos 2k\theta=\sum_{k=1}^n\sin 2k\theta=0\), so \(\sum\cos^2k\theta=\sum\sin^2k\theta=n/2\) and \(\sum\cos k\theta\sin k\theta=0\). The answer is \(1-\dfrac{n^2}{4}\).

\problem{5}
\hint{1}
\(\operatorname{adj}A\) is nilpotent, hence singular, hence by the rank facts about the adjugate its rank is \(1\) or \(0\).
\hint{2}
A rank-one matrix is \(uv^T\), and \((uv^T)^m=(v^Tu)^{m-1}uv^T\). So nilpotency forces \(v^Tu=0\), and then already \((uv^T)^2=(v^Tu)uv^T=O\).

\problem{6}
\hint{1}
Look for a product that the two relations make vanish. Try computing \((A^2+B^2)(A-B)\).
\hint{2}
\((A^2+B^2)(A-B)=A^3-A^2B+B^2A-B^3=0\). Since \(A-B\ne 0\), \(A^2+B^2\) cannot be invertible. The answer is no.

\problem{7}
\hint{1}
By Theorem 3, the hypothesis says \(\det(A+kB)=\pm1\) for \(k=0,\dots,4\). Consider \(f(t)=\det(A+tB)\): what kind of function of \(t\) is it?
\hint{2}
\(f\) is a polynomial of degree at most \(2\) with integer coefficients taking values in \(\{1,-1\}\) at five points. By pigeonhole it takes the same value at three of them, so it is constant. Hence \(\det(A+5B)=\pm1\).

\problem{8}
\hint{1}
It suffices to show the kernel is trivial (dimension count). Suppose \(\varphi(A)=0\) with \(A\ne0\). What does determinant-preservation give for \(\det(X+A)\) for arbitrary \(X\)?
\hint{2}
\(\det(X+A)=\det(\varphi(X)+\varphi(A))=\det\varphi(X)=\det X\) for all \(X\). Now find an \(X\) that is singular but with \(X+A\) invertible.
\hint{3}
If \(\operatorname{rank}A=r\ge1\), write \(A=P\operatorname{diag}(E_r,0)Q\) with \(P,Q\) invertible and take \(X=P\operatorname{diag}(0,E_{n-r})Q\). Then \(\det X=0\) but \(X+A=PQ\) is invertible, a contradiction.
\end{document}
